\documentclass[12pt,a4j]{jarticle}
\usepackage{graphicx}
\begin{document}
\title{基礎プログラミングおよび演習 レポート ＃06}
\author{2120029 政野玄空}
\date{2021年11月14日}
\maketitle

\section{構想・計画・設計}

美しいと思う絵を生成するということで美しいと思うものをまず考えた。結果猫を描くことに決定した。猫の全体像を描くのは非常に難易度が高いので顔面だけを描くことにした。
その後猫のサンプル写真として知り合いの飼っているノルウェージャンフォレストキャットの写真を選んだ。美しい。

まず顔面を中心に描くこととしてど真ん中の座標、耳を置く位置、鼻、口、ひげの大まかな座標を決めた。
その後実装しながら細かい調整を行い完成した。
\section{プログラムコード}
\begin{verbatim}
  Pixel = Struct.new(:r, :g, :b)
  $img = Array.new(200) do
    Array.new(300) do Pixel.new(255,255,255) end
  end
  def pset(x, y, r=0, g=0, b=0, a=0.0)
    if x < 0 || x >= 300 || y < 0 || y >= 200 then return end
    $img[y][x].r = ($img[y][x].r * a + r * (1.0 - a)).to_i
    $img[y][x].g = ($img[y][x].g * a + g * (1.0 - a)).to_i
    $img[y][x].b = ($img[y][x].b * a + b * (1.0 - a)).to_i
  end
  def writeimage(name)
    open(name, "wb") do |f|
      f.puts("P6\n300 200\n255")
      $img.each do |a|
        a.each do |p| f.write(p.to_a.pack("ccc")) end
      end
    end
  end
  def fillcircle(x, y, rad, r=0, g=0, b=0, a=0.0)
    j0 = (y-rad).to_i; j1 = (y+rad).to_i
    i0 = (x-rad).to_i; i1 = (x+rad).to_i
    j0.step(j1) do |j|
      i0.step(i1) do |i|
        if (i-x)**2+(j-y)**2<rad**2
          if block_given? then yield(i,j) else pset(i,j,r,g,b,a) end
        end
      end
    end
  end
  def filldonut(x, y, r1, r2, r=0, g=0, b=0, a=0.0)
    j0 = (y-r1).to_i; j1 = (y+r1).to_i
    i0 = (x-r1).to_i; i1 = (x+r1).to_i
    j0.step(j1) do |j|
      i0.step(i1) do |i|
        d2 = (i-x)**2+(j-y)**2
        if r2**2 <= d2 && d2 <= r1**2
          if block_given? then yield(i,j) else pset(i,j,r,g,b,a) end
        end
      end
    end
  end
  def fillrect(x, y, w, h, r=0, g=0, b=0, a=0.0)
    j0 = (y-0.5*h).to_i; j1 = (y+0.5*h).to_i
    i0 = (x-0.5*w).to_i; i1 = (x+0.5*w).to_i
    j0.step(j1) do |j|
      i0.step(i1) do |i|
        if block_given? then yield(i,j) else pset(i,j,r,g,b,a) end
      end
    end
  end
  def fillellipse(x, y, rx, ry, r=0, g=0, b=0, a=0.0)
    j0 = (y-ry).to_i; j1 = (y+ry).to_i
    i0 = (x-rx).to_i; i1 = (x+rx).to_i
    j0.step(j1) do |j|
      i0.step(i1) do |i|
        if ((i-x).to_f/rx)**2 + ((j-y).to_f/ry)**2 < 1.0
          if block_given? then yield(i,j) else pset(i,j,r,g,b,a) end
        end
      end
    end
  end
  def fillrotellipse(x, y, rx, ry, theta, r=0, g=0, b=0, a=0.0)
    d = (if rx > ry then rx else ry end)
    j0 = (y-d).to_i; j1 = (y+d).to_i
    i0 = (x-d).to_i; i1 = (x+d).to_i
    j0.step(j1) do |j|
      i0.step(i1) do |i|
        dx = i - x; dy = j - y;
        px = dx*Math.cos(theta) - dy*Math.sin(theta)
        py = dx*Math.sin(theta) + dy*Math.cos(theta)
        if (px/rx)**2 + (py/ry)**2 < 1.0
          if block_given? then yield(i,j) else pset(i,j,r,g,b,a) end
        end
      end
    end
  end
  def filltriangle(x0, y0, x1, y1, x2, y2, r = 0, g = 0, b = 0, a = 0.0)
    if block_given?
      fillconvex([x0, x1, x2, x0], [y0, y1, y2, y0]) do |x,y| yield(x,y) end
      fillconvex([x0, x2, x1, x0], [y0, y2, y1, y0]) do |x,y| yield(x,y) end
    else
      fillconvex([x0, x1, x2, x0], [y0, y1, y2, y0], r, g, b, a)
      fillconvex([x0, x2, x1, x0], [y0, y2, y1, y0], r, g, b, a)
    end
  end
  def fillconvex(ax, ay, r=0, g=0, b=0, a=0.0)
    xmax = ax.max.to_i; xmin = ax.min.to_i
    ymax = ay.max.to_i; ymin = ay.min.to_i
    ymin.step(ymax) do |j|
      xmin.step(xmax) do |i|
       if isinside(i, j, ax, ay)
         if block_given? then yield(i,j) else pset(i,j,r,g,b,a) end
       end
      end
    end
  end
  def isinside(x, y, ax, ay)
    (ax.length-1).times do |i|
      if oprod(ax[i+1]-ax[i],ay[i+1]-ay[i],x-ax[i],y-ay[i])<0
        return false
      end
    end
    return true
  end
  def oprod(a, b, c, d)
    return a*d - b*c;
  end
  def fillline(x0, y0, x1, y1, w, r=0, g=0, b=0, a=0.0)
    dx = y1-y0; dy = x0-x1; n = 0.5*w / Math.sqrt(dx**2 + dy**2)
    dx = dx * n; dy = dy * n
    if block_given?
      fillconvex([x0-dx, x0+dx, x1+dx, x1-dx, x0-dx],
                 [y0-dy, y0+dy, y1+dy, y1-dy, y0-dy]) do |x,y| yield(x,y) end
    else
      fillconvex([x0-dx, x0+dx, x1+dx, x1-dx, x0-dx],
                 [y0-dy, y0+dy, y1+dy, y1-dy, y0-dy], r, g, b, a)
    end
  end
  
  fillellipse(150,100,132,82)
  fillellipse(150,100,130,80,255,255,255)
  fillellipse(150,60,100,40,245,222,179)
  filltriangle(79, 3, 59, 61, 122, 61)
  filltriangle(214, 3, 189, 61, 242, 61)
  filltriangle(80, 5, 60, 60, 120, 60,255,182,193)
  filltriangle(215, 5, 190, 60, 240, 60,255,182,193)
  fillrotellipse(90, 100, 10, 20, -30)
  fillrotellipse(210, 100, 10, 20, 30)
  filltriangle(130, 120, 170, 120, 150, 150,255,182,193)
  filltriangle(150, 150, 110, 160, 180, 160,255,182,193)
  fillline(10, 120, 100, 145, 1)
  fillline(10, 150, 100, 150, 1)
  fillline(10, 180, 100, 155, 1)
  fillline(200, 145, 290, 120, 1)
  fillline(200, 150, 290, 150, 1)
  fillline(200, 155, 290, 180, 1)
  writeimage("t.ppm")
  
\end{verbatim}

\section{プログラムの説明}
サンプルからメソッドをコピーして持ってくる。
その後まず顔の輪郭を描写する。
その後顔の全体を描写する。
次に耳となる三角を生成する。輪郭も同時にコピペして少し外側にでるように調整する。
次に目玉を2つ描画する。330度で調整すると思ったとおりにならず、マイナスを利用したらいい感じになった。
次に鼻を描画する。
次に口を描画する。
最後にひげを6本薄い長方形を利用し描画する。
いままでの命令をt.ppmに書き込む。
(プログラムのどの部分が何をしているかを説明する)

\section{生成された絵}
猫の顔面の絵。参考にした絵の美しさとは程遠いが猫に見えるので間違いなく美しいと言える。

\begin{center}
\includegraphics[width=12cm]{cat.ps}
\end{center}


\section{考察}
最初に大まかな設計をし、色は何でもいいのでわかるようにわけて画像を生成しながら微調整を重ねていったのが功を奏して短い時間で完成できたと考えられる。ただ先に同等のドット絵などを準備していればもっと早く完成できた可能性がある。

\section{アンケート}

\subsection{Q1：画像が自由に生成できるようになりましたか。}
できた。

\subsection{Q2：画像をうまく生成する「コツ」は何だと思いましたか。}
先にドット絵を生成するツールなどをつかったりして写真を取り込んでおけば座標がかんたんに取れると思った。

\subsection{Q3：リフレクション(今回の課題で分かったこと)・感想・要望をどうぞ。}
猫の美しさを表現するのはとても困難だった。

\end{document}